% ===== §18 =====
\section{The Guarding of Generativity}
\label{p26:sec:justicegen}

\noindent
This section takes the third register of the limit's generativity, justice, and establishes that justice at its root is the guarding of a generativity it cannot settle, and that the mode of this guarding is a knowing that holds what it cannot dispose of, which the paper names witnessing. The section takes up the limit set out in Part One, that justice owes protection to a generativity it cannot bring into the settled sign, gives the positive account of justice that follows from it, reviews the work on witnessing that fixes the mode, and draws the claim.

\subsection*{Beneath settlement}

Justice is commonly figured as settlement, the bringing of a matter to a close by a judgment that disposes of it. The account of this paper places another office beneath that one. Part One established that what a relation of justice most owes protection to is the generativity of relational being, the still-forming life that the intervention of the Real keeps opening in the order of settled meaning, and that this generativity cannot be settled in advance, since to settle it is to convert the still-becoming into the already-become and so to end the very thing that was to be protected. Justice, at its root, is therefore not settlement but guarding: the holding-open of a space in which generativity can go on, against the pressure to close it prematurely into a finished case. Settlement remains a real and necessary office of justice, addressed to what has taken form and can be brought before a judgment; but beneath it lies the deeper office of keeping open what has not yet taken form, and it is this that the limit grounds.

\subsection*{The mode of guarding, witnessing}

The mode of this guarding is a knowing that does not settle, and the paper names it, again as a conclusion of its own path, witnessing. To witness is to hold what one cannot dispose of, to be turned toward a generativity or a suffering that one does not bring to a close by pronouncing on it. The work on testimony fixes the mode precisely. Felman and Laub found that the witnessing of what exceeds settlement is not the recording of a fact but a holding of what has no adequate statement, a keeping-faith with what cannot be closed into a case, such that the witness anchors the thing without substituting a judgment for it \parencite{felman1992testimony}. Justice, insofar as it guards a generativity it cannot settle, must know in this mode. It must hold the still-forming and the unstated without pretending to have disposed of them, anchoring them in the common life without closing them into the settled sign. A justice built on witnessing does not mistake the settled sign for the whole of what is owed; it disposes of what can be disposed of and keeps open what cannot.

\begin{claim}[Justice as witnessing]
The root office of justice is to guard a generativity it cannot settle, and the mode of that guarding is witnessing: a knowing that holds the still-forming and the unstated without disposing of them, anchoring them in the common life without closing them into the settled sign. Justice as witnessing keeps its own procedures answerable to what they have not foreseen, and refuses itself the false completeness of a settlement that has anticipated everything.
\end{claim}

\subsection*{Against an administration of the already-formed}

This is the just face of the limit's generativity, and its distance from a justice of pure settlement is exact. A justice that could settle everything in advance would be an administration of the already-formed, closed to the new that relational being keeps generating and deaf to the suffering that no procedure exhausts. The justice the limit grounds keeps itself open to what it cannot settle, guards the generativity it cannot close, and finds in the unreachability of that generativity not a gap in the law but the reason the law exists. The three registers of Part Four now stand together, and their unity is visible. The poem gives circulation without capture; ethics gives a bond without possession; justice gives a guarding without settlement. In each, the same refusal of closure that Part Two showed to be the condition of the drive's life is shown to be the condition of a generative form: in language, in the relation to the other, and in the common life. What cannot be brought into the sign is, in every register, not a failure but the opening in which the generative form becomes possible.
